% ============================================================
% Báo cáo mục 4: Trao đổi khóa Diffie-Hellman và ứng dụng thực tế
% Người viết: Tô Tuệ Nghi. Nguồn tham khảo nằm trong references.bib (BibTeX).
% Dòng "% KIEM CHUNG" là chỗ phải đối chiếu lại trước khi nộp.
% ============================================================

\section{Trao đổi khóa Diffie-Hellman và ứng dụng thực tế}
\label{sec:dh-ungdung}

\subsection{Giao thức và liên hệ với chương trình}

Giao thức Diffie-Hellman \cite{diffie1976new} cho phép hai bên chưa có bí mật chung
thiết lập một khóa chung thông qua một kênh công khai. Độ an toàn dựa trên bài
toán logarit rời rạc \cite{slide-bai2}. Hai bên dùng chung số nguyên tố lớn $p$
và phần tử sinh $g$ của nhóm $\mathbb{Z}_p^*$; các bước được tóm tắt trong
Bảng~\ref{tab:dh-steps}.

\begin{table}[!htbp]
\centering
\begin{tabular}{@{}lll@{}}
\toprule
Bước & Alice & Bob \\
\midrule
1. Chọn khóa riêng & chọn ngẫu nhiên $x$ & chọn ngẫu nhiên $y$ \\
2. Tính giá trị công khai & $k_A \equiv g^x \pmod p$ & $k_B \equiv g^y \pmod p$ \\
3. Trao đổi & gửi $k_A$ cho Bob & gửi $k_B$ cho Alice \\
4. Tính bí mật chung & $k \equiv k_B^{\,x} \pmod p$ & $k' \equiv k_A^{\,y} \pmod p$ \\
\bottomrule
\end{tabular}
\caption{Các bước của giao thức Diffie-Hellman theo \cite{slide-bai2}.}
\label{tab:dh-steps}
\end{table}

Sơ đồ trong Hình~\ref{fig:dh-flow} mô tả cùng quá trình theo trình tự thời gian.

\begin{figure}[!htbp]
\centering
\begin{tikzpicture}[>=latex, font=\small]
  % Hai bên
  \node[draw, rounded corners, minimum width=3cm, minimum height=0.8cm] (A) at (0,0) {Alice};
  \node[draw, rounded corners, minimum width=3cm, minimum height=0.8cm] (B) at (10,0) {Bob};
  \draw[dashed] (0,-0.4) -- (0,-9.4);
  \draw[dashed] (10,-0.4) -- (10,-9.4);

  % Tham số công khai
  \node[draw, fill=gray!10, rounded corners, align=center, text width=8cm]
       at (5,-1.5) {Tham số công khai: số nguyên tố an toàn $p$ và phần tử sinh $g$};

  % Khóa riêng
  \node[draw, fill=white, rounded corners, align=center, text width=3.9cm]
       at (0,-3.1) {Chọn khóa riêng $x$\\(giữ bí mật)};
  \node[draw, fill=white, rounded corners, align=center, text width=3.9cm]
       at (10,-3.1) {Chọn khóa riêng $y$\\(giữ bí mật)};

  % Giá trị công khai
  \node[draw, fill=white, rounded corners, align=center, text width=3.9cm]
       at (0,-4.7) {$k_A \equiv g^x \pmod p$};
  \node[draw, fill=white, rounded corners, align=center, text width=3.9cm]
       at (10,-4.7) {$k_B \equiv g^y \pmod p$};

  % Trao đổi qua kênh công khai
  \draw[->, thick] (0,-5.8) -- node[above] {$k_A$} (10,-5.8);
  \draw[<-, thick] (0,-6.6) -- node[above] {$k_B$} (10,-6.6);
  \node[align=center, text width=5.4cm, font=\footnotesize\itshape] at (5,-7.35)
       {kênh công khai: kẻ nghe lén chỉ thấy $p$, $g$, $k_A$, $k_B$};

  % Bí mật chung
  \node[draw, fill=white, rounded corners, align=center, text width=3.9cm]
       at (0,-8.3) {$k \equiv k_B^{\,x} \pmod p$};
  \node[draw, fill=white, rounded corners, align=center, text width=3.9cm]
       at (10,-8.3) {$k' \equiv k_A^{\,y} \pmod p$};
  \node[draw, fill=gray!10, rounded corners, align=center, text width=6cm]
       at (5,-9.9) {$k = k' \equiv g^{xy} \pmod p$ (bí mật chung)};
\end{tikzpicture}
\caption{Sơ đồ trao đổi khóa Diffie-Hellman giữa Alice và Bob.}
\label{fig:dh-flow}
\end{figure}

Hai bên thu được cùng một giá trị vì
\[
k \equiv k_B^{\,x} \equiv (g^y)^x \equiv g^{xy} \equiv (g^x)^y \equiv k_A^{\,y} \equiv k' \pmod p.
\]
Kẻ nghe lén biết $p$, $g$, $k_A$, $k_B$ nhưng muốn tính $g^{xy}$ thì phải giải
bài toán logarit rời rạc; theo tài liệu môn học, phương pháp
index-calculus được xem là mạnh nhất hiện biết cho bài toán này \cite{slide-bai2}.

Chương trình của nhóm cài đúng các bước trên: sinh số nguyên tố an toàn $p$
(\texttt{generate\_\allowbreak safe\_\allowbreak prime}), chọn phần tử sinh $g$
(\texttt{find\_\allowbreak generator}), sinh khóa riêng trong đoạn $[2,\,p-2]$
(\texttt{generate\_\allowbreak private\_\allowbreak key}), rồi dùng \texttt{modular\_\allowbreak exponentiation} để
tính $k_A$, $k_B$ và hai bí mật chung. Khóa riêng được chọn trong $[2,\,p-2]$ vì với $x \in \{0,\,1,\,p-1\}$ thì
$g^x \bmod p$ bằng $1$ hoặc bằng chính $g$, làm lộ khóa.
% DA KIEM CHUNG 11/11/2026 voi main.cpp va dh.cpp tren main: ten ham va thu tu cac buoc dung.
% DA KIEM CHUNG 11/11/2026 voi random.cpp tren main: generate_private_key tra ve khoa rieng trong [2, p-2] (random_range lay mau loai bo, khoang dong).

\subsection{Ứng dụng chính: thiết lập khóa phiên trong TLS 1.3}

TLS là giao thức bảo vệ HTTPS. Phiên bản 1.3 hiện được đặc tả trong RFC~9846
(tháng 7/2026), thay thế RFC~8446 (tháng 8/2018); bản mới là một cập nhật nhỏ,
giữ nguyên số phiên bản và tương thích ngược \cite{rfc9846}. Giao thức hỗ trợ ba
chế độ thiết lập khóa: trao đổi khóa bất đối xứng, khóa chia sẻ trước (PSK) và
kết hợp cả hai. Trao đổi khóa bất đối xứng gồm Diffie-Hellman trên trường hữu hạn
hoặc trên đường cong elliptic \cite{rfc9846}. Chương trình của nhóm tương ứng với
trường hợp trên trường hữu hạn $\mathbb{Z}_p^*$.

Cách hoạt động khi dùng trao đổi khóa bất đối xứng \cite{rfc9846}: trong thông
điệp \texttt{ClientHello}, máy khách gửi phần chia sẻ khóa
(\texttt{key\_\allowbreak share}) của mình; máy chủ trả lời trong \texttt{ServerHello} bằng
phần chia sẻ \emph{tạm thời} của máy chủ, cùng nhóm với một phần chia sẻ của máy
khách. Từ hai phần chia sẻ này, hai bên cùng tính được bí mật chung, rồi dùng nó
để dẫn xuất khóa bảo vệ dữ liệu. Với nhóm hữu hạn, giá trị công khai có dạng
$g^{x} \bmod p$ (RFC~9846 viết là $Y = g^X \bmod p$, với $X$ là khóa riêng), và mỗi bên \emph{phải} kiểm tra giá trị công khai nhận được
thỏa $1 < h < p-1$ (với $h$ là giá trị công khai của đối phương) để đối phương không ép hệ thống vào một nhóm con nhỏ
\cite{rfc9846}.

\paragraph{Khóa tạm thời và bảo mật về sau (forward secrecy).}
TLS 1.3 đã loại bỏ các bộ mã dùng RSA tĩnh và Diffie-Hellman tĩnh; mọi cơ chế
thiết lập khóa dựa trên khóa công khai đều có forward secrecy \cite{rfc9846}.
Nghĩa là nếu khóa dài hạn của máy chủ bị lộ sau này, các phiên đã ghi lại trước
đó vẫn không giải mã được, miễn là hai bên hủy khóa tạm thời ngay khi dùng xong
\cite{rfc7919}. RFC~9846 còn cấm dùng lại một phần chia sẻ khóa cho nhiều kết nối
\cite{rfc9846}. Vì vậy ở chế độ mặc định, chương trình của nhóm lấy seed từ hệ điều hành và sinh
khóa riêng mới cho mỗi lần chạy thay vì dùng khóa cố định; tham số \texttt{--seed}
dùng seed cố định và chỉ phục vụ kiểm thử.
% DA KIEM CHUNG 11/11/2026 voi main.cpp tren main (seed_from_os la che do mac dinh).

\paragraph{Các nhóm chuẩn.}
TLS 1.3 đã bỏ các nhóm Diffie-Hellman tùy chọn, chỉ dùng các nhóm có tên; với
trường hữu hạn đó là \texttt{ffdhe2048} đến \texttt{ffdhe8192}, định nghĩa trong
RFC~7919 \cite{rfc9846,rfc7919}. Các nhóm này dùng số nguyên tố an toàn
\cite{rfc7919}; độ mạnh ước tính nêu ở Bảng~\ref{tab:ffdhe}.
% DA KIEM CHUNG 11/11/2026 voi RFC 7919 Phu luc A: 103/125/150/175/192 bit dung; g = 2 cho ca 5 nhom.

\begin{table}[!htbp]
\centering
\begin{tabular}{@{}lcc@{}}
\toprule
Nhóm & Kích thước $p$ (bit) & Độ mạnh (bit) \\
\midrule
\texttt{ffdhe2048} & 2048 & 103 \\
\texttt{ffdhe3072} & 3072 & 125 \\
\texttt{ffdhe4096} & 4096 & 150 \\
\texttt{ffdhe6144} & 6144 & 175 \\
\texttt{ffdhe8192} & 8192 & 192 \\
\bottomrule
\end{tabular}
\caption{Các nhóm Diffie-Hellman hữu hạn trong RFC~7919 \cite{rfc7919}. Độ mạnh là số bit khóa đối xứng có độ an toàn tương đương.}
\label{tab:ffdhe}
\end{table}

\subsection{Tấn công man-in-the-middle và vai trò của xác thực}

Giao thức ở Bảng~\ref{tab:dh-steps} không có bước xác thực danh tính: $k_A$ và
$k_B$ được gửi đi nguyên dạng. Một kẻ tấn công chủ động nằm giữa đường truyền có
thể giả làm Bob với Alice và giả làm Alice với Bob, tự thiết lập một khóa chung
riêng với mỗi bên, rồi đọc và sửa mọi dữ liệu chuyển tiếp. Đây là phân tích của
nhóm dựa trên mô tả giao thức, không phải kết luận trích từ một nguồn cụ thể.

TLS chặn tấn công này bằng cách kết hợp trao đổi khóa với xác thực. Máy chủ gửi
chứng chỉ, một chữ ký trên toàn bộ quá trình bắt tay (\texttt{CertificateVerify})
và một mã xác thực (\texttt{Finished}), nhờ đó định danh của máy chủ được gắn với
khóa vừa thỏa thuận \cite{rfc9846}. Cơ chế chữ ký số này chính là nội dung của
bài \emph{Chữ ký số} trong tài liệu môn học \cite{slide-bai4}: người gửi ký bằng
khóa cá nhân, người nhận xác minh bằng khóa công khai.

\subsection{Vì sao cần số nguyên tố an toàn}

Số nguyên tố $p$ là \emph{số nguyên tố an toàn} nếu $q = (p-1)/2$ cũng nguyên tố,
khi đó $p-1 = 2q$. Có hai lợi ích cụ thể.

\paragraph{Tìm phần tử sinh dễ dàng.}
Thuật toán tìm phần tử sinh trong tài liệu môn học cần các ước nguyên tố của
$p-1$ và thử $g^{(p-1)/p_i} \bmod p$ với từng ước $p_i$ \cite{slide-bai2}. Với số
nguyên tố an toàn, chỉ có hai ước là $2$ và $q$, nên chỉ cần kiểm tra
$g^2 \not\equiv 1$ và $g^q \not\equiv 1 \pmod p$. Nếu $p$ bất kỳ thì phải phân
tích $p-1$ ra thừa số nguyên tố, việc này khó khi $p$ lớn. Hàm
\texttt{find\_\allowbreak generator} của nhóm làm đúng như vậy: chọn $g$ ngẫu nhiên,
kiểm tra hai điều kiện trên và lặp lại đến khi thỏa. Vì $g$ được chọn trong $[2,\,p-2]$ và $p$ nguyên tố nên $g^2 \equiv 1 \pmod p$ chỉ xảy ra khi $g \equiv \pm 1$, tức $g \in \{1,\,p-1\}$; do đó điều kiện thứ nhất luôn đúng và thực chất chỉ còn phải kiểm tra $g^q \not\equiv 1$. Chương trình vẫn kiểm cả hai điều kiện cho khớp với phân tích bậc của $g$ (bằng $1$, $2$, $q$ hoặc $2q$); chi phí thêm là một phép lũy thừa $g^2$, không đáng kể.
% DA KIEM CHUNG 11/11/2026 voi dh.cpp tren main (q = (p-1)>>1, kiem tra g^2 != 1 va g^q != 1).
Vì nhóm có $\varphi(p-1)=q-1$
phần tử sinh trong tổng số $p-1$ phần tử, mỗi lần thử một $g$ ngẫu nhiên thành
công với xác suất xấp xỉ $1/2$.

\paragraph{Kiểm tra giá trị công khai đơn giản.}
Khi biết nhóm dùng số nguyên tố an toàn, nhóm có kích thước $p-1$ hoặc $(p-1)/2$,
và chỉ cần kiểm tra $1 < Y < p-1$ để biết đối phương không ép hệ thống vào nhóm
con hai phần tử \cite{rfc9846,rfc7919}. Ngược lại, với nhóm tùy chọn do máy chủ gửi, máy
khách không biết cấu trúc của nhóm, và nếu thấy nhóm không đủ mạnh ($p$ quá nhỏ,
$p$ không nguyên tố, hoặc có nhóm con nhỏ khó tránh) thì chỉ còn cách từ chối kết
nối \cite{rfc7919}.

\subsection{Kích thước tham số trong thực tế}

Đề bài yêu cầu $p$ tối thiểu 512 bit; đây là mức dành cho bài tập, không phải mức
dùng trong thực tế. Nghiên cứu Logjam cho thấy sau một tuần tính trước cho một
nhóm 512 bit cố định, có thể tính logarit rời rạc bất kỳ trong nhóm đó chỉ khoảng
một phút; 82\% máy chủ dễ tổn thương dùng cùng một nhóm 512 bit, và các trình
duyệt lớn đã đổi sang từ chối nhóm ngắn \cite{logjam-cacm}. Nguyên nhân sâu xa là
kẻ tấn công chỉ cần tính trước một lần cho một số nguyên tố $p$ rồi dùng lại cho
mọi mục tiêu dùng chung $p$ đó \cite{logjam-cacm}. RFC~7919 yêu cầu máy khách phải
từ chối nhóm tùy chọn có $|p| < 768$ bit và nên từ chối nhóm dưới 1024 bit; với
hệ thống tương lai, theo RFC~7919, ENISA khuyến nghị nhóm từ 3072 bit trở lên \cite{rfc7919}.
% DA KIEM CHUNG 11/11/2026: nguong 768/1024 bit o RFC 7919 muc 3.1; khuyen nghi ENISA 3072 bit o muc 8.2.

\subsection{Các ứng dụng khác}

Diffie-Hellman còn là cơ chế trao đổi khóa trong hai giao thức quen thuộc khác.
\emph{SSH}: RFC~4253 mô tả trao đổi khóa Diffie-Hellman (mục 8) và định nghĩa các
phương thức \texttt{diffie-hellman-\allowbreak group1-\allowbreak sha1} và
\texttt{diffie-hellman-\allowbreak group14-\allowbreak sha1} \cite{rfc4253}; RFC~8268 bổ sung các phương
thức từ \texttt{diffie-hellman-\allowbreak group14-\allowbreak sha256} đến
\texttt{diffie-hellman-\allowbreak group18-\allowbreak sha512} với nhóm lớn hơn và hàm băm SHA-2
\cite{rfc8268}. \emph{IPsec}: trong giao thức IKEv2, lần trao đổi đầu tiên
(\texttt{IKE\_\allowbreak SA\_\allowbreak INIT}) thương lượng tham số bảo mật, gửi số dùng một lần và gửi
các giá trị Diffie-Hellman \cite{rfc7296}; IKEv2 dùng DH hoặc ECDH để thiết lập
bí mật chung \cite{rfc9370}. Nghiên cứu Logjam cũng chỉ ra rằng một số ít nhóm cố
định hoặc chuẩn hóa được dùng cho hàng triệu máy chủ: tính trước cho một nhóm
1024 bit đã đủ để nghe lén thụ động 18\% trang HTTPS phổ biến, và một nhóm thứ hai
đủ để giải mã lưu lượng của 66\% IPsec VPN và 26\% máy chủ SSH
\cite{logjam-cacm}.
% DA KIEM CHUNG 11/11/2026 voi tom tat bai bao Logjam (CACM 2019): 82%, mot tuan/mot phut, 18%, 66%, 26%.
